% ---
% title: Three exams in a row
% tags: multiplication-rule, conditional-probability
% ---

\begin{problem}
	A graduate takes three professional (actuarial) exams, one after the other.
	After a failure she stops and takes no further exam. She passes the first
	exam with probability $0.9$. If she passed the first exam, she passes the
	second with probability $0.8$. If she passed the first two exams, she passes
	the third with probability $0.7$.
	(a) What is the probability that she passes all three exams?
	(b) What is the probability that she fails exactly at the third exam?
	(c) She did not pass all three exams. What is the conditional probability
	that she failed the second exam?
\end{problem}

\begin{slideproblem}
	Three exams in a row; after a failure she stops. $P(\text{pass 1st})=0.9$;
	if she passed the 1st: $P(\text{pass 2nd})=0.8$; if she passed the first
	two: $P(\text{pass 3rd})=0.7$. Probability that she (a) passes all three;
	(b) fails exactly at the third; (c) failed the second, given that she did
	not pass all three?
\end{slideproblem}

\begin{slides}
	\small
	\exstatement
	\begin{columns}[T]
		\begin{column}{.5\textwidth}
			\uncover<2->{\sol $E_i$ = ``she passes exam $i$''.
				The numbers are \emph{conditional}:\\[.5mm]
				$P(E_1)=0.9$,\quad $P(E_2\mid E_1)=0.8$,\\
				$P(E_3\mid E_1E_2)=0.7$.}\\[1mm]
			\uncover<3->{Tree: these numbers on the branches.}\\[1mm]
			\uncover<4->{Multiplication rule (multiply along the path):
				\[ P(E_1E_2E_3)=P(E_1)\,P(E_2\mid E_1)\,P(E_3\mid E_1E_2). \]}%
			\uncover<5->{(a) $P(E_1E_2E_3)=0.9\cdot0.8\cdot0.7=\alert{0.504}$.}
		\end{column}
		\begin{column}{.48\textwidth}
			\centering
			\uncover<3->{%
				\begin{tikzpicture}[scale=.78,
						lab/.style={font=\scriptsize},
						ev/.style={font=\scriptsize,inner sep=1pt}]
					\coordinate (S) at (0,0);
					\node[ev] (P1) at (1.6,.9) {$E_1$};
					\node[ev] (F1) at (1.6,-.9) {$E_1^c$};
					\node[ev] (P2) at (3.2,1.6) {$E_2$};
					\node[ev] (F2) at (3.2,.3) {$E_2^c$};
					\node[ev] (P3) at (4.8,2.1) {$E_3$};
					\node[ev] (F3) at (4.8,1.1) {$E_3^c$};
					\fill (S) circle (1.5pt);
					\draw[navy,thick] (S) -- node[lab,above left=-1pt] {0.9} (P1);
					\draw[gold,thick] (S) -- node[lab,below left=-1pt] {0.1} (F1);
					\draw[navy,thick] (P1) -- node[lab,above left=-1pt] {0.8} (P2);
					\draw[gold,thick] (P1) -- node[lab,below left=-1pt] {0.2} (F2);
					\draw[navy,thick] (P2) -- node[lab,above left=-1pt] {0.7} (P3);
					\draw[gold,thick] (P2) -- node[lab,below left=-1pt] {0.3} (F3);
					\node[lab,right=1pt] at (F1.east) {stops};
					\node[lab,right=1pt] at (F2.east) {stops};
					\uncover<5->{\node[lab,right=1pt,navy] at (P3.east)
						{$0.504$};}
				\end{tikzpicture}}
		\end{column}
	\end{columns}
\end{slides}

\begin{slides}
	\small
	\begin{columns}[T]
		\begin{column}{.52\textwidth}
			(b) Fails exactly at the third: passes 1st, passes 2nd, fails 3rd:
			\[ P(E_1E_2E_3^c)=0.9\cdot0.8\cdot0.3=\alert{0.216}. \]
			\uncover<2->{(c) Condition: \emph{not all three},
				$(E_1E_2E_3)^c$. Complement rule:
				\[ P\big((E_1E_2E_3)^c\big)=1-0.504=0.496. \]}%
			\uncover<3->{Failed the 2nd: $P(E_1E_2^c)=0.9\cdot0.2=0.18$.
				Check: $0.1+0.18+0.216=0.496$.}\\[1mm]
			\uncover<4->{Failing the 2nd means she did not pass all three: the
				intersection is $E_1E_2^c$ itself.}
		\end{column}
		\begin{column}{.46\textwidth}
			\centering
			\begin{tikzpicture}[scale=.78,
					lab/.style={font=\scriptsize},
					ev/.style={font=\scriptsize,inner sep=1pt}]
				\coordinate (S) at (0,0);
				\node[ev] (P1) at (1.6,.9) {$E_1$};
				\node[ev] (F1) at (1.6,-.9) {$E_1^c$};
				\node[ev] (P2) at (3.2,1.6) {$E_2$};
				\node[ev] (F2) at (3.2,.3) {$E_2^c$};
				\node[ev] (P3) at (4.8,2.1) {$E_3$};
				\node[ev] (F3) at (4.8,1.1) {$E_3^c$};
				\fill (S) circle (1.5pt);
				\draw[navy,thick] (S) -- node[lab,above left=-1pt] {0.9} (P1);
				\draw[gold,thick] (S) -- node[lab,below left=-1pt] {0.1} (F1);
				\draw[navy,thick] (P1) -- node[lab,above left=-1pt] {0.8} (P2);
				\draw[gold,thick] (P1) -- node[lab,below left=-1pt] {0.2} (F2);
				\draw[navy,thick] (P2) -- node[lab,above left=-1pt] {0.7} (P3);
				\draw[gold,thick] (P2) -- node[lab,below left=-1pt] {0.3} (F3);
				\node[lab,right=1pt] at (P3.east) {$0.504$};
				\node[lab,right=1pt,teal] at (F3.east) {$0.216$};
				\uncover<3->{\node[lab,right=1pt,teal] at (F2.east)
					{$0.18$};}
				\uncover<3->{\node[lab,right=1pt,teal] at (F1.east)
					{$0.1$};}
			\end{tikzpicture}
		\end{column}
	\end{columns}
	\uncover<5->{\[ P\big(E_1E_2^c\mid (E_1E_2E_3)^c\big)
			=\frac{P(E_1E_2^c)}{P\big((E_1E_2E_3)^c\big)}
			=\frac{0.18}{0.496}=\frac{45}{124}\approx\alert{0.363}. \]}
\end{slides}

\begin{solution}
	\idea The exams come one after the other, and each probability in the
	statement is \emph{conditional} on what happened before. A probability tree
	shows this: one branching per exam, the conditional probabilities on the
	branches. The probability of a path is the product of the probabilities
	along it (the multiplication rule).

	\step{Events and data.} Let $E_i$ = ``she passes exam $i$'' ($i=1,2,3$).
	The statement gives
	\[ P(E_1)=0.9,\qquad P(E_2\mid E_1)=0.8,\qquad P(E_3\mid E_1E_2)=0.7. \]
	The second number is \emph{not} $P(E_2)$: it is the probability of passing
	the second exam among the cases where she passed the first. By the complement
	rule for a fixed condition, the failure probabilities are
	\[ P(E_1^c)=0.1,\qquad P(E_2^c\mid E_1)=0.2,\qquad
		P(E_3^c\mid E_1E_2)=0.3. \]

	\begin{center}
		\begin{tikzpicture}[scale=1.05,
				lab/.style={font=\small},
				ev/.style={inner sep=1.5pt}]
			\coordinate (S) at (0,0);
			\node[ev] (P1) at (2,1.1) {$E_1$};
			\node[ev] (F1) at (2,-1.1) {$E_1^c$};
			\node[ev] (P2) at (4,1.9) {$E_2$};
			\node[ev] (F2) at (4,.4) {$E_2^c$};
			\node[ev] (P3) at (6,2.5) {$E_3$};
			\node[ev] (F3) at (6,1.3) {$E_3^c$};
			\fill (S) circle (1.5pt);
			\draw[navy,thick] (S) -- node[lab,above left=-1pt] {0.9} (P1);
			\draw[gold,thick] (S) -- node[lab,below left=-1pt] {0.1} (F1);
			\draw[navy,thick] (P1) -- node[lab,above left=-1pt] {0.8} (P2);
			\draw[gold,thick] (P1) -- node[lab,below left=-1pt] {0.2} (F2);
			\draw[navy,thick] (P2) -- node[lab,above left=-1pt] {0.7} (P3);
			\draw[gold,thick] (P2) -- node[lab,below left=-1pt] {0.3} (F3);
			\node[lab,right=2pt] at (P3.east) {all three: $0.504$};
			\node[lab,right=2pt] at (F3.east) {fails 3rd: $0.216$};
			\node[lab,right=2pt] at (F2.east) {fails 2nd: $0.18$};
			\node[lab,right=2pt] at (F1.east) {fails 1st: $0.1$};
		\end{tikzpicture}
	\end{center}
	After a failure she stops, so a failure branch ends the path. The tree has
	four paths, and every possible history is exactly one of them.

	\step{(a)} Passing all three is the top path, $E_1E_2E_3$. The
	multiplication rule says: the probability that several events all happen is
	the probability of the first, times the probability of the second given the
	first, times the probability of the third given the first two:
	\[ P(E_1E_2E_3)=P(E_1)\,P(E_2\mid E_1)\,P(E_3\mid E_1E_2)
		=0.9\cdot0.8\cdot0.7=0.504. \]
	In words: 90\% pass the first exam; 80\% of these pass the second
	($0.72$); 70\% of those pass the third ($0.504$).

	\step{(b)} Failing exactly at the third exam means passing the first two
	and then failing the third: the path $E_1E_2E_3^c$. Multiply along it:
	\[ P(E_1E_2E_3^c)=P(E_1)\,P(E_2\mid E_1)\,P(E_3^c\mid E_1E_2)
		=0.9\cdot0.8\cdot0.3=0.216. \]

	\step{(c)} The condition is ``she did not pass all three'', the complement
	of $E_1E_2E_3$. By the complement rule,
	\[ P\big((E_1E_2E_3)^c\big)=1-0.504=0.496. \]
	Check with the tree: not passing all three means failing the first, the
	second or the third exam (three separate paths, which cannot happen
	together):
	\[ 0.1+0.18+0.216=0.496. \]
	``She failed the second exam'' means she passed the first and failed the
	second: the path $E_1E_2^c$, with probability
	\[ P(E_1E_2^c)=P(E_1)\,P(E_2^c\mid E_1)=0.9\cdot0.2=0.18. \]
	If she failed the second exam, she certainly did not pass all three, so
	$E_1E_2^c$ lies inside $(E_1E_2E_3)^c$ and their intersection is
	$E_1E_2^c$ itself. The definition of conditional probability gives
	\[ P\big(E_1E_2^c\mid(E_1E_2E_3)^c\big)
		=\frac{P(E_1E_2^c)}{P\big((E_1E_2E_3)^c\big)}
		=\frac{0.18}{0.496}=\frac{45}{124}\approx0.363. \]
	\begin{answer}
		(a) $0.504$ \quad (b) $0.216$ \quad (c) $\frac{45}{124}\approx0.363$
	\end{answer}
	\begin{remark}
		\textbf{Common mistake:} in (c), writing the event as $E_2^c$ alone.
		She also does not pass the second exam if she fails the first one and
		stops, so $P(E_2^c)=0.1+0.18=0.28$. ``Failed the second exam'' means she
		took it and failed: $E_1E_2^c$.
	\end{remark}
	\begin{remark}
		Among the graduates who did not pass all three exams, $0.1/0.496\approx
		0.202$ failed the first, $0.363$ the second and $0.216/0.496\approx0.435$
		the third. These three conditional probabilities add up to $1$: given the
		condition, exactly one of the three failure paths happened.
	\end{remark}
\end{solution}
